% Folder 83, batch 06 (archivists' pages 101-120).
% Pass: Opus 5.5 (claude-opus-5-5), 2026-10-10 — first pass, unchecked against the pages by a human.
%
\documentclass[11pt,a4paper]{article}
\input{../preamble/grothendieck}

\folder{83}
\batch{6}
\pages{101}{120}
\foldertitle{[Icosaèdres, bi-icosaèdres et algèbres d'Azumaya de rang 4] : notes manuscrites (1982-1983, 1986, s.d.).}
\dating{1982-1986}
\watermark{Édition de démonstration}

\begin{document}

\section{Azumaya rang 4 (II)}
\note{Titre de sa main en tête de la p. 101, souligné.}

\page{101}

Avril 86
\note{Date écrite en haut à droite de la p. 101.}

Soit $A$ Alg.\ d'Azumaya rang 4 sur \struck{schéma $S$} \add{anneau $k$ (disons)}, et $B\subset A$ un sous-\struck{Module loc.\ libre de} \add{fibré vectoriel de} rang 3. Alors $B^{\perp}$ est un sous-\struck{vectoriel} \add{fibré vectoriel} de rang 1, soit $D$. On a corr.\ 1-1 entre les $B$, et les $D$, car la forme quadratique canonique $\operatorname{Tr}(uv)$ sur $A$ est non dégénérée.

À vrai dire, il y a deux formes bil.\ sym.\ non dégénérées intéressantes sur $A$, savoir
\[
\varphi(u,v)=\operatorname{Tr}(uv) \quad\text{et}\quad \psi(u,v)=\varphi_{\det}(u,v)=\operatorname{Tr}u\operatorname{Tr}v-\operatorname{Tr}(uv),
\]
\note{« $\psi(u,v)=$ » est ajouté au-dessus de $\varphi_{\det}$.}
\struck{donnant ayant des discriminants qui \ill{}}

Je dis que pour le fibré vectoriel $B$ (de rang quelconque) qui contient $1$, \emph{les deux formes \uncertain{induites} sur $B$ \uncertain{ont même} discriminant}. En effet, il suffit de le voir dans un cas universel d'un tel $B$ sur $\operatorname{Spec}\mathbb{Z}$, ce qui nous permet de supposer $2$ inversible, donc
\[
A\simeq k.1\oplus A^{\natural}, \qquad B=k.1\oplus V,\quad V\subset A^{\natural},
\]
\[
A^{\natural}\overset{\mathrm{def}}{=}\{u\in A \mid \operatorname{Tr}u=0\} = 1^{\perp}
\]
(\emph{NB} L'orthogonal de $1$ \uncertain{pour l'une ou l'autre forme est le même}\,!)\note{L'exposant de $A^{\natural}$ est un signe en forme d'éclair (ou de « 4 » couché) ; il est rendu partout par $\natural$.}

\page{102}
et \struck{$\xi$} on est ramené à prouver l'assertion \struck{pour} analogue pour $k.1$ et $V$ (régularité du disc.). Or sur $k1$ les deux formes \struck{bil.} \uncertain{quadr.\ ass.}\ prennent la \uncertain{valeur} $2$, OK. Sur $V$, \struck{c'est} elles sont \uncertain{égales}\,! On gagne\,!

\struck{Proposition}

\emph{NB} Soit $B$ sous-module de $A$ tel que $B\ni 1$, alors
\[
B^{\perp(\varphi)}=B^{\perp(\psi)} .
\]
En effet, \struck{les deux $B^{\perp}$ \ill{} orthogonaux c'est une $B^{\perp}$ est \ill{} de} on a
\[
\varphi(1,u)=\psi(1,u)=\operatorname{Tr}u
\]
donc les orthogonaux de $1$ sont les mêmes (c'est $A^{\natural}=\{u\in A\mid \operatorname{Tr}u=0\}$), et $B^{\perp}\subset A^{\natural}$. Or si $\xi\in A^{\natural}$, on a : pour \struck{$\varphi$ \ill{}} $u\in A$
\[
\psi(u,\xi)=-\varphi(u,\xi)
\]
donc dire que $\xi$ est orth.\ à $B$ pour $\varphi$, ou pour $\psi$, est pareil, cqfd.

\page{103}
Ça peut être plus explicité encore. \struck{Plaçons} Les $\delta_{B}$ pour $B$ variable forment une section du fibré $\underline{\mathcal{O}}_{P}(2)$ \add{$=\underline{\mathcal{O}}_{P}(1)^{\otimes 2}$} sur le fibré projectif (de dimension relative 2)
\[
\mathbb{P}(\dot{A}/k)
\]
des (hyper)plans dans $A$, donc provenant d'un élément canonique de
\[
\Gamma(\underline{\mathcal{O}}_{P}(2))\simeq \operatorname{Sym}^{2}_{k}(A/k)\simeq \text{espace des formes quadratiques sur } (A/k)^{\vee}\simeq A^{\natural} .
\]
Or il y a une telle forme canonique sur $A^{\natural}$, savoir $\det$. Et on vérifie que $\delta_{B}$ est bien défini par cet élément. (Cela n'a rien à voir avec la structure multiplicative de $A$, qui intervient pour définir les $\delta_{B}$, il suffit de la forme quadratique \add{non dég.} $\det$, et l'élément $1$ de $A$ \struck{\ill{} injection}…)\note{Le point sur le $A$ de $\mathbb{P}(\dot{A}/k)$ est lu tel quel ; lecture incertaine.}

\page{104}
\note{La p. 104 commence au milieu d'un raisonnement. L'ordre de lecture semble être, à rebours de la numérotation des archivistes : p. 108 (Proposition), p. 107 (Dém.), p. 106, p. 105, p. 104, chaque feuillet continuant la phrase laissée en suspens au bas du précédent.}
\[
\mu:\underbrace{\textstyle\bigwedge^{2}B/k}_{=\ \det B/k}\longrightarrow A/B
\]
Or on a un iso.\ can.
\[
\det B\simeq \underbrace{\det k}_{k}\otimes\det B/k\simeq \det B/k
\]
\marginal{attention à convention de signes}
donc on peut interpréter $\mu$ comme un
\[
\mu:\det B\longrightarrow A/B
\]
\note{Ici comme dans quelques autres formules hors ligne de ce lot (p. 105, 106, 114, 117), un court fragment biffé et illisible a été omis de la formule.}
avec $\det B\simeq (A/B)^{\otimes -1}$, soit encore
\[
\mu\in A/B^{\otimes 2}\simeq D^{-\otimes 2}.
\]
Ceci dit, il est clair que
\[
B \text{ sous-algèbre} \iff \mu=0 .
\]
Il suffit de voir que
\[
\mu=\pm\delta
\]
(le signe dépendant des conventions de signes). Or il suffit de le voir dans le cas universel $S=\mathbb{P}^{2}_{\mathbb{Z}}$, et pour \uncertain{ceux-ci} \struck{\ill{}} $\underline{\mu}$ et $\underline{\delta}$ \struck{\ill{}} universels, \uncertain{et} les \ill{} \struck{zéros, \ill{}} ont des zéros, -- gagné\,!

\page{105}
laquelle disparaît quand on passe au carré :
\[
(\det B)^{\otimes 2}\otimes (A/B)^{\otimes 2}\simeq \underline{1}
\]
qui donne
\[
(\det B)^{\otimes -2}\simeq (A/B)^{\otimes 2}.
\]
Ainsi, on peut \uncertain{ici voir}
\[
\delta_{\det B}\in (A/B)^{\otimes 2}\simeq D^{\otimes -2}
\]
[et en fait, c'est simplement donné par la restriction de la forme quadratique $\det$ : $D$ : $\delta_{\det B}=\det|_{D}$]

Or on a une application bilin.\ can.
\[
\begin{aligned}
B\times B &\longrightarrow A/B\\
(u,v)&\longmapsto uv \bmod B
\end{aligned}
\]
et celle-ci est alternée, puisque pour H.C., on a pour $u\in B$
\[
u^{2}\equiv 0 \bmod B
\]
(car $u^{2}=(\operatorname{Tr}u)u-\det u.1$, et $1,u\in B$). Ainsi, on trouve
\[
\textstyle\bigwedge^{2}B\xrightarrow{\ \mu\ }A/B .
\]
On sait que
\[
\mu(u\wedge v)=uv \bmod B .
\]
D'ailleurs \struck{\ill{}} la forme bil.\ $uv \bmod B$ est nulle quand l'un ou l'autre des \ill{} est dans $k$, donc on \struck{peut} considère $\mu$ comme\note{Après la deuxième occurrence de « $uv \bmod$ », un signe ressemblant à « $1\cdot$ » précède le $B$ ; lecture incertaine.}

\page{106}
correspondant
\[
B_{D}=D^{\perp}
\]
(= normalisateur de $D$ dans l'alg.\ de Lie $A$ i.e.\ ens.\ des $x\in A$ tels que \struck{$[u,x]\in D$ pour \ill{} $x\in D$} $\operatorname{ad}_{\mathrm{Lie}}(x)(D)\subset D$)

($B_{D}^{*}$ = normalisateur de $D$ dans $A^{\times}$)

D'où équivalence de (i) avec (i'). D'autre part, (ii') $\Longrightarrow$ (ii) trivialement, il reste à voir que (iii) $\Longrightarrow$ (i) -- je ne suis pas si sûr que c'est vrai, (\uncertain{sauf} sur un corps (ou un \uncertain{anneau} \ill{} réduit) le \uncertain{cas} \ill{} décisif est celui où $k=k_{0}[t]/t^{2}$ … (l'histoire des \ill{} d'une algèbre de Borel…)\note{L'argument sur (i), (i'), (ii), (ii'), (iii) renvoie à des énoncés qui ne figurent pas dans ce lot.}

\struck{Conclusion}. Notons que $\delta_{B}=\operatorname{disc}'_{\det B}$ est une section de
\[
\delta_{B}\in(\det B)^{\otimes -2}\Bigl[\underset{\text{iso can}}{\simeq} A/B^{\otimes 2}\Bigr]
\]
compte tenu que l'on a un iso can
\[
\det(A)\simeq k
\]
\[
\det B\otimes A/B\simeq\det A\ (\simeq k)
\]
Cet iso dépend d'une convention de signes,

\page{107}
\emph{Dém.} L'équivalence de (ii) et (ii') est évidente, du fait que puisque $A$ a une forme quadratique (vg quelc.) non dégénérée, et qu'on regarde \struck{les} \add{des} \struck{\ill{}} sous-fibrés vectoriels de codim 1 (\uncertain{on suppose} \ill{} pair) \struck{\ill{}} si $2$ inv.\ dans $k$ -- \struck{quand \ill{} dans ce} dernier cas, \ill{} de polar.\ du disc.\ \emph{divisé}, le discriminant \uncertain{ordinaire} (\ill{} $(A)$),
\[
\operatorname{disc}_{\varphi_{B}}=\operatorname{disc}_{\varphi_{D}}=2\operatorname{disc}'_{\det}(B)
\]
\note{L'indice du deuxième discriminant est lu $\varphi_{D}$ sans certitude.}
\struck{\ill{}} \uncertain{en vertu du} disc., la condition (ii) \uncertain{implique} (ii'), et la réciproque est vraie si $2$ inv.

\struck{\ill{}} Si $D$ est isotrope, il est orthogonal à lui-même, donc
\[
D\subset D^{\perp}=B .
\]
L'inverse est vrai si $2$ est inv.\ (\emph{NB} si $2=0$ dans $k$, donc on a \emph{tjrs} $D\subset B$\,!) En effet, $D\subset B$ signifie que $D$ est orthogonal à lui-même, et donc
\[
2d(x)=\varphi_{d}(x,x)=0 .
\]
Cela \uncertain{veut dire} que $2d(x)=0$ \uncertain{ou} $d(x)=0$ pour $x\in D$ i.e.\ $D$ est isotrope.

On sait que les sous-alg.\ de Borel sont en corr.\ 1-1 avec les \struck{\ill{}} $D$ isotropes \add{dans $A^{\natural}$}, à \struck{\ill{}} une telle $D$

\page{108}
Les sous-fibrés vectoriels $B$ de $A$, (de rang 3, \add{contenant 1,}) correspondent par orthogonalité (pour $\varphi$ ou $\psi$, c'est kif kif) aux sous-fibrés vectoriels de rang 1 $D$ de $A^{\natural}$.

D'autre part, pour un tel $B$, considérons
\[
\det_{B}=\det|B ,
\]
forme quadratique sur $B$, elle a donc un discriminant divisé
\[
\operatorname{disc}'(\det_{B}) \quad\text{tel que}\quad 2\operatorname{disc}'(\det_{B})=\operatorname{disc}\psi_{B}
\]
où $\psi_{B}$ est la restriction de $\psi$ à $B$. Ainsi
\[
\operatorname{disc}\psi_{B}=\operatorname{disc}\varphi_{B}=2\underbrace{\operatorname{disc}'(\det_{B})}_{\overset{\mathrm{def}}{=}\ \delta(B)}
\]

\emph{Proposition} Les \add{trois} conditions suivantes sont équivalentes
\begin{itemize}
\item[(i)] $D$ est \emph{isotrope} i.e.\ $\det|D=0$
\item[(ii)] $\delta(B)=0$ \quad (\emph{NB} $\delta(B)\in$ \struck{\ill{}} $\det(B)^{\otimes(-2)}$)
\item[\struck{\ill{}}(iii)] $B$ est une sous-algèbre de $A$
\item[(iii')] \struck{\ill{}} $B$ est une « sous-alg.\ de Borel » de $A$
\end{itemize}
\marginal{\struck{Si $S$ est le spectre d'un corps, \ill{} (iii) \ill{}}}\note{Note marginale à gauche, barrée, reliée par une flèche à (iii).}
Ces \struck{\ill{}} conditions impliquent les deux suivantes, et si $2.1_{k}$ est \uncertain{inv.}, \struck{elles} \uncertain{en sont} équivalentes aux précédentes
\begin{itemize}
\item[(iv)] $D\subset B$
\item[(ii')] \struck{disc} $\operatorname{disc}(\varphi_{B})=0$ \quad ($\Longleftrightarrow \operatorname{disc}_{\psi_{B}}=0$)
\end{itemize}

\page{109}
\emph{Cor.} Conditions équivalentes
\begin{itemize}
\item[(i)] $D$ est anisotrope i.e.\ $\det_{A}|D$ est non nulle sur chaque fibre
\item[(ii)] $\delta(B)$ ($=\delta(D)$) est un élément inv.\ de $(\det B)^{\otimes -2}$ \struck{\ill{}}

\struck{i.e.\ $\det_{A}|D$ est partout non nulle}

i.e.\ $\det_{B}=\det_{A}|B$ est non dégénérée
\item[(iii)] $B$ n'est une sous-alg.\ sur aucune fibre
\item[(iv)] $B$ \add{fibre} engendre $A$ comme \emph{algèbre} (pour \uncertain{chaque} \uncertain{fibre} \uncertain{bien sûr})
\item[(v)] $D\not\subset B$ sur les fibres
\item[(vi)] $\varphi|B$ est non dégénérée
\end{itemize}
\marginal{si $2$ inv.\ sur $S$}\note{La note marginale, reliée par une accolade, porte sur (v) et (vi).}

En effet, si $u,v$ sont deux éléments de $B$ tels que $1,u,v$ forment une base, les conditions équivalent à :
\begin{itemize}
\item[(v)] $1,u,v,uv$ forment une base de $A$
\item[(v')] $uv \bmod B$ est un \struck{\ill{}} \add{élément} inversible de $A/B$
\end{itemize}

\struck{Supposons}

Nous regardons maintenant deux sous-\struck{\ill{}} \add{fibrés} \struck{\ill{}} $H,H'$ de rang 2, contenant 1 (ce sont \uncertain{nécess.} des sous-algèbres, par H.C.) et supposons
\[
H\cap H'=k
\]
donc
\[
H\oplus H'\overset{\mathrm{déf}}{=}B \quad\text{sur la fibre sous-fibré de rang 3}
\]
\note{Lecture de la dernière ligne incertaine : « sur la fibre » semble se rattacher à la condition $H\cap H'=k$.}

\page{110}
\struck{\ill{}}, Un $H$ \add{$\ni 1$} est défini aussi par son orthogonal, et de même pour $H'$
\[
V\subset A^{\natural},\quad V'\subset A^{\natural}
\]
sous-fibrés vectoriels de rang 2, qui satisfont donc
\[
V+V'=A^{\natural}
\]
donc
\[
V\cap V'\overset{\mathrm{def}}{=}D
\]
est un sous-fibré vect.\ de rang 1.

\struck{Dire que $H$ est \ill{} que $H$ est}

Les $H$ correspondent \add{sur $k$ ou $k$-points} \struck{\ill{}} aux sections \add{des} \struck{\ill{}} \add{plans} du \struck{\ill{}} \add{projectif} $\check{\mathbb{P}}(A/k)$, et les $V$ aux droites dans le plan projectif dual, dans lequel il y a une conique \uncertain{canonique}, celle des droites isotropes dans $A^{\natural}$. Alors $H$ est \uncertain{déployé} sur $k$ (\uncertain{les} $H$ qui \uncertain{agissent sur} $k$) \struck{\ill{}} \add{géométriquement} ssi \uncertain{elle} \uncertain{définit} \uncertain{un} \struck{point} fibre, la droite (plaine) $D_{H}=D_{V}$ est \struck{\ill{}} tangente à la conique, i.e.\ \ill{} découpe une diviseur \struck{\ill{}} non réduit : un pt. Si cette condition est vérifiée pour

\marginal{\emph{NB} Les $H$ \ill{} correspondent 1-1 aux \ill{} \ill{} \ill{} \ill{} \ill{} \ill{} corr.\ diviseur \ill{} \ill{}}\note{Note marginale oblique en bas à gauche, en grande partie illisible ; elle semble porter sur les $H$ en correspondance 1-1 avec des diviseurs (« diviseurs de degré 2 » ?).}

\page{111}
$H$ et $H'$, on dit qu'ils sont \emph{en position générale} si sur chaque fibre géométrique, les deux paires de points sur la conique isotrope sont disjointes. Mais cela signifie qu'il n'y a pas de \struck{\ill{}} droites isotropes (sur la fibre) orth.\ à la \emph{fois} à $H$ et $H'$, i.e.\ à $B$. Donc par la \uncertain{même} proposition, cela signifie que $\delta_{B}$ est inv.\ \uncertain{ici}, que $\det|D$ est inversible. \struck{\ill{}} Dans ce cas, \struck{élucidé} p.\ ex.\ le choix (possible localement) de \struck{$H,H'$ \ill{}} \struck{Spin \ill{}} \add{droites isotropes} \uncertain{disjointes}
\[
D_{1},D_{2}\subset V\cap \underset{\substack{\parallel\\ \text{con.\ isotrope}}}{C},\qquad D_{1}',D_{2}'\subset V'\cap C
\]
\note{Le $V$ de $V\cap C$ est écrit sur un autre symbole.}
on a un \uncertain{rapport} anharmonique
\[
\lambda=\lambda_{(D_{1},D_{2})}(D_{1}',D_{2}')\in k \qquad \lambda \text{ partout } \neq 0,1 \text{ i.e. } \lambda \text{ et } 1-\lambda \text{ inv.}
\]
qui se transforme en $1/\lambda$ si on échange $D_{1}$ et $D_{2}$ (\ill{} $D_{1}'$ et $D_{2}'$), et qui se \uncertain{dispose}

\page{112}
pas de l'ordre dans chaque paire \uncertain{prise} les deux \struck{paires} couples $(D_{1},D_{2})$, $(D_{1}',D_{2}')$.

On a un \uncertain{rapport} \uncertain{plus} intrinsèque, qui ne dépend pas de permutations
\[
D_{1}\longleftrightarrow D_{2} \qquad\text{ou}\qquad D_{1}'\longleftrightarrow D_{2}'
\]
\[
\rho=\lambda+\lambda^{-1}
\]
qui est défini sur $k$ (sans localiser). On retrouve $\lambda$, au signe près : $\lambda^{-1}$ puis, comme une solution de
\[
\lambda^{2}-\lambda\rho+1=0 .
\]
\struck{Le $\rho$ est \uncertain{pour}} On a ici
\[
\rho\neq 2 \quad\text{sur chaque fibre,}
\]
i.e.
\[
\rho_{0}=\rho-2 \quad\text{est }\neq 0\text{ sur chaque fibre i.e.\ est inv.}
\]
On voit \struck{que} aisément que les classes de \struck{paires} \struck{dans} $(H,H')$ \add{(H, H' étales sur $k$, $H,H'$ en pos.\ gén.)} dans $A_{4}$ \add{modulo conjugaison local} correspondent aux diff.\ $\rho_{0}$ inv.\ i.e.\ si deux telles paires $(H,H')$ avec inv.\ $\rho_{0}=\rho_{0}(H,H')$, $(\overline{H},\overline{H}')$ avec inv.\ $\overline{\rho_{0}}$,

\page{113}
alors $\exists$ \add{localt} \add{$\Theta$} automorphisme \struck{\ill{}} de $A$, tel que
\[
\Theta(H)=\overline{H} \qquad \Theta(H')=\overline{H'}
\]
ssi
\[
\overline{\rho_{0}}=\rho_{0}
\]
D'autre part, si on se donne $\rho_{0}$, \add{(\uncertain{inv.})} il existe \emph{localt} une paire $(H,H')$ qui le réalise.

\struck{D'ailleurs, \ill{} il existe un automorphisme de $A$ \ill{} \ill{} \ill{} loc.\ \ill{} ($\mathrm{Pic}(k)=0$), et pour $H$ il existe loc.\ \ill{} \add{propre} loc.\ \ill{}, il fixe $H$ (il est \ill{}… propre auto.\ \ill{} dans $A$) ou \ill{} provient d'un $\Theta\in$ \ill{} de $k$ \ill{} \ill{} $x\mapsto x'=\operatorname{Tr}x\,1-x$ \ill{} \ill{} certain \uncertain{posant}}\note{Passage barré par trois longs traits obliques et par des traits horizontaux ; seuls quelques mots se lisent.}

Dans $G=\underline{\mathrm{Aut}}(A)\simeq\underline{\mathrm{Aut}}(P)$, le groupe qui stabilise $H$ et $H'$ est un sous-\struck{gr} schéma \ill{} isom.\ : $(\mathbb{Z}/2\mathbb{Z})_{S}$, \struck{\ill{}} dont \ill{} \ill{} est un \ill{} \ill{} \ill{} $H$ et sur $H'$ \struck{\ill{}}

\marginal{Si $\rho\neq -2$ partout, \struck{\ill{}} \ill{} $\rho_{0}\neq -4$ partout, \ill{} \ill{} $\simeq\mathbb{Z}/2\mathbb{Z}$ \ill{} $\overline{C}$ \ill{} \ill{}}\note{Note marginale oblique en bas à gauche, partiellement lisible.}

\page{114}
induit sur $H$ et $H'$ l'automorphisme
\[
x\longmapsto x'=\operatorname{Tr}(x).1-x
\]
En tout cas, les isom.\ qui envoient $(H,H')$ sur $(\overline{H},\overline{H}')$ forment un torseur sous le groupe $\mathbb{Z}/2\mathbb{Z}$ (si $\rho_{0}+4$ inv.) chaque torseur peut (je pense \ill{}) être un trivial (\ill{} \ill{}) un \uncertain{le} torseur associé du groupe
\[
\underline{\mathrm{Aut}}(A)\simeq \check{V}(A)^{*}/\mathbb{G}_{m} \quad\text{est trivial}
\]

\struck{\ill{}} Écrit plutôt : [Si $P$ est un \add{$\mathbb{Z}/2$-}torseur sur $\operatorname{Spec}k$, on a un \uncertain{homom.}
\[
P\xrightarrow{\ \sim\ }\operatorname{Isom}(\varepsilon_{H},\varepsilon_{\overline{H}})\ (=\varepsilon_{H}\wedge\varepsilon_{\overline{H}})
\]
\[
\text{\struck{$P\times\varepsilon_{H}\longrightarrow\varepsilon_{\overline{H}}$}}
\]
\struck{qui définit} et donc
\[
P\xrightarrow{\ \sim\ }\varepsilon_{H'}\wedge\varepsilon_{\overline{H}'} \quad ]
\]
\marginal{$\varepsilon_{H}$ désigne le revêt.\ étale de $S=\operatorname{Spec}k$ défini par $H$}

\struck{\ill{}} $P$ \uncertain{muni} de sections \add{$\varepsilon_{H'}\simeq\varepsilon_{\overline{H}'}$} pas si $\varepsilon_{H}\simeq\varepsilon_{\overline{H}}$ ; \uncertain{prenant} par exemple
\[
H'=\overline{H}' \quad\text{correspondant : la bisection}
\]
$\{0,\infty\}$ de $\mathbb{P}^{1}$, et $H,\overline{H}$ étant les bisections définies resp.\ par
\[
\begin{aligned}
Z^{2}-bZ+c&=0 \qquad b^{2}-4c\neq 0\\
Z^{2}-\overline{b}Z+\overline{c}&=0 \qquad \overline{b}^{2}-4\overline{c}\neq 0
\end{aligned}
\]
\marginal{sur un corps…}

\page{115}
on a ici
\[
\rho=\frac{b^{2}}{c}-2,\qquad \overline{\rho}=\frac{\overline{b}^{2}}{\overline{c}}-2
\]
la condition $\rho=\overline{\rho}$ s'écrit
\[
\frac{b^{2}}{c}=\frac{\overline{b}^{2}}{\overline{c}},
\]
qui n'implique pas l'isom.\ des deux extensions\,! P.\ ex.\ en car.\ $\neq 2$, si on prend $b=\overline{b}=0$ la condition \struck{$(\rho=-2)$} (il est vrai que cas exceptionnel ($\rho=-2$)) est satisfaite \uncertain{par} \uncertain{l'existence} dans le cas d'isom.\ entre les \uncertain{deux}, mais les \uncertain{\ill{}} \uncertain{ne} \uncertain{sont} \uncertain{pas} \uncertain{isom.} \struck{\ill{}} $\varepsilon_{H}$ et $\varepsilon_{\overline{H}}$ quand si $\overline{c}/c$ est un non carré. \struck{\ill{}} Etc.…

Il serait nécessaire de voir si tant $\rho_{0}$ inv.\ peut être réalisé loc.\ \emph{Zar} tel un seulement \uncertain{\ill{}} (\ill{}) \note{passage peu sûr.}

Je suppose que oui ; un point d'un $H$, et un pour un $H'=$ \uncertain{$\operatorname{int}(u)H$} pour $u\in A^{*}$ convenable…

\emph{NB} Vérification immédiate dans le cas où $A$ est scindée…

\page{116}
Je vais donner une expression de $\rho$. D'abord on trouve d'une base
\[
1,u,v \qquad\text{base de } B=H+H'
\]
posant
\[
\begin{cases}
b=\operatorname{Tr}u,\ c=\det u & \delta_{H}=b^{2}-4c\\
b'=\operatorname{Tr}v,\ c'=\det v & \delta_{H'}=b'^{2}-4c'\\
t=\operatorname{Tr}uv=\operatorname{Tr}vu
\end{cases}
\]
\emph{NB} La condition $H$ étale $\Leftrightarrow\delta_{H}$ inv., $H'$ étale $\Leftrightarrow\delta_{H'}$ inv.

on trouve l'expression
\[
\boxed{\rho_{0}=\frac{P}{Q}}
\]
où
\[
\begin{aligned}
P&=P(b,c,b',c',t)=(b^{2}-4c)(b'^{2}-4c')=\delta_{H}\delta_{H'}\\
Q&=Q(b,c,b',c',t)=t^{2}-bb't+(b^{2}c'+b'^{2}c-4cc')
\end{aligned}
\]
\note{Après la parenthèse de $Q$, un signe raturé en marge droite.}
et d'autre part on a
\[
Q=\tfrac{1}{2}\delta_{\varphi|B} \text{ pour la base } 1,u,v
=\tfrac{1}{2}\det\begin{pmatrix}\operatorname{Tr}1&\operatorname{Tr}u&\operatorname{Tr}v\\ \operatorname{Tr}u&\operatorname{Tr}u^{2}&\operatorname{Tr}uv\\ \operatorname{Tr}v&\operatorname{Tr}vu&\operatorname{Tr}v^{2}\end{pmatrix}
\]
\[
=\text{discr.\ divisé de la forme } \det_{|B} \text{ p.r.\ à la base } 1,u,v
\]
\note{Devant $\tfrac12$ et devant « discr. » un signe (« $-$ » ?) est noirci.}
(matrice \add{de $\varphi|B$} égale :
\[
\begin{pmatrix}2&b&b'\\ b&b^{2}-2c&t\\ b'&t&b'^{2}-2c'\end{pmatrix}\ )
\]
\note{Le coefficient $b^{2}-2c$ est écrit sur un premier chiffre raturé.}
donc par hyp.\ ($H,H'$ en position générale) \struck{donc} on a bien $Q$ inversible, donc $\rho_{0}$ est défini, et non nul.

Je voudrais donner : cette formule pour $\rho_{0}$ une forme intrinsèque, indépendante du choix de bases.

\page{117}
Regardons d'abord un sous-fibré \struck{vect.} de rg.\ 2 quelconque, $H$, et regardons \struck{\ill{}} la forme \struck{det} quadratique
\[
\det|H
\]
Donc on peut considérer son discriminant \add{$\delta_{H}$}, est aussi \uncertain{égal} à celui de la forme bilinéaire \add{sym.} $\varphi|H$ \uncertain{($\operatorname{Tr}uv$)} (car $H\ni 1$). C'est une section de
\[
\det H^{\otimes(-2)}\simeq (H/k)^{\otimes -2}
\]
Par rapport à une base $1,u$, d'où la base $1\wedge u$ de $\det H\simeq H/k$, et $(1\wedge u)^{\otimes -2}$ de $(\det H)^{\otimes -2}$, on trouve le discriminant
\[
\det\begin{pmatrix}\operatorname{Tr}1&\operatorname{Tr}u\\ \operatorname{Tr}u&\operatorname{Tr}u^{2}\end{pmatrix}=b^{2}-4c
\]
donc
\[
\delta_{H}=(b^{2}-4c)(1\wedge u)^{\otimes -2}
\]
Si on a $H,H'$ tels que $H\cap H'=k$ (sous-coch.\ de rang 2 en position transversale \uncertain{p.r.} $k.1$), et posant
\[
B=H+H'
\]
comme dessus, on a
\[
B/k\simeq H/k+H'/k
\]
d'où isom.\ (défini à un signe près, si on précise les conventions)
\[
\det B\simeq \underbrace{\det k}_{k}\otimes\det B/k\simeq H/k\otimes H'/k
\]
Donc \struck{\ill{}}
\[
(\det B)^{\otimes -2}\simeq (H/k)^{\otimes -2}\otimes(H'/k)^{\otimes -2}
\]

\page{118}
isom.\ qui ne dépend plus des signes. \struck{Ceci} On peut donc considérer
\[
\begin{cases}
\delta_{H}\delta_{H'}\in(\det B)^{\otimes -2}\simeq D^{\otimes -2}\\
\delta_{B}\in(\det B)^{\otimes -2}\simeq D^{\otimes -2}
\end{cases}
\]
Je rappelle que la \struck{\ill{}} condition que \struck{\ill{}} $H,H'$ soient « en position générale » i.e.\ engendrent l'\emph{algèbre} $A$, est que $\delta_{B}$ est section \emph{inversible} de $(\det B)^{\otimes -2}\simeq D^{\otimes -2}$ i.e.\ $D$ anisotrope. Ceci dit, on
\[
\boxed{\rho_{0}=\frac{\delta_{H}\delta_{H'}}{\delta_{B}}}
\]
\marginal{trouve : si $H,H'$ ét.\ \uncertain{sur} $k$ i.e.\ $\delta_{H}\delta_{H'}$ section inversible}
Cette formule n'est définie que si $H,H'$ sont étales $/S$ et en position générale ($\delta_{B}$ \add{inv.}). Mais on peut la prendre comme \emph{définition} de $\rho_{0}$ dès que $\delta_{B}$ inv.\ i.e.\ $H,H'$ en position générale (pas nécess.\ étales). On peut définir $\rho_{0}$ encore plus généralement, dès que $\delta_{H}\delta_{H'}$ et $\delta_{B}$ ne

\page{119}
s'annulent pas simultanément, i.e.\ ils engendrent $D^{\otimes -2}$ : on trouve alors pour $\rho_{0}$ une section de $\mathbb{P}^{1}_{k}$ prenant la valeur $\infty$ aux pts où $\delta_{B}$ s'annule, \struck{la valeur} i.e.\ $H,H'$ contiennent une \add{(mais $H,H'$ étales)} \emph{Borel}, et la valeur $0$ aux les pts où $H$ ou $H'$ n'est pas étale.

\emph{NB} Le seul cas sur un corps où $\rho_{0}$ n'est pas défini, est celui où $B=B_{D}$ est une sous-alg.\ de Borel, pour droite $D\subset$ \struck{\ill{}} $B$ isotr.\ ) et $H$ (disons, i.e.\ un \uncertain{des} \uncertain{deux} rôles) est $k\oplus D$, et enfin $H'$ est un \add{plan de $B$} \struck{\ill{}} qui contient $k$ et est distinct de $H$ \add{sans plus}. \struck{Cette situation} Deux telles situations sont conjuguées loc.\ \ill{} i.e.\ ce type de situation est \uncertain{fixé} au choix d'une Borel seulement [correspond maximal \uncertain{\ill{}} drapeau, i.e.\ Borel et d'un tore i.e.\ \emph{couple de Killing}.]
\note{Fin de page : la phrase entre crochets se lit mal ; « maximal » et « drapeau » sont peu sûrs.}

\page{120}
Revenons à l'expression explicite
\[
-\delta_{B=H+H'}=Q=t^{2}-bb't+(b^{2}c'+b'^{2}c-4cc')
\]
\note{Devant le premier et le second signe « $=$ », des lettres raturées.}
Je dis que pour un syst.\ de 5 \uncertain{constantes} $(b,c,b',c',t)\in k^{5}$ tels que cette quantité soit inversible*, il existe toujours, \add{une forme} \struck{loc.\ \ill{} \ill{} $\operatorname{Spec}k$,} des les \uncertain{réalisations} dans une algèbre d'Azumaya $A$ \uncertain{par} \uncertain{deux} éléments $u,v$, \struck{Pour les} tels que $H=k\oplus ku$ et $H'=k\oplus kv$ soient des sous-fibrés vect.\ de rg 2 en position relative générale. En fait \uncertain{on} \uncertain{connaît} déjà $(H,u)$ \struck{\ill{}} comme \uncertain{donnée} à \uncertain{isom.\ près} \uncertain{près} par $b,c$, et \uncertain{itou} pour $(H',v')$ \uncertain{par} $b',c'$. Si $(b^{2}-4c)(b'^{2}-4c')$ \add{$\delta\delta'$} est inversible, on trouve une quantité $\rho_{0}\neq 0$, qui correspond à \struck{deux} une position relative dans $\mathbb{M}_{2\,k}$ disons \uncertain{à} \uncertain{condition} de faire extension de la base \uncertain{vers} les solutions de
\[
\lambda^{2}-(\rho_{0}+2)\lambda+1=0
\]
\marginal{* Sans doute il suffit que $\delta$ soit inversible ; on fait une \uncertain{vérif.} \uncertain{plus} \uncertain{loin}…}
\marginal{À partir d'ici, \uncertain{nous} \uncertain{supposons} $\delta$ \uncertain{inversible}, \uncertain{sans} \ill{} \ill{} \ill{} \ill{}}
\note{Deux notes marginales obliques à gauche, la seconde en grande partie illisible. L'argument se poursuit au-delà de ce lot.}

\end{document}
